A query follows rotation $q'=\Pi q$, IVF routing, primary screening and residual refinement. Final selection returns $k$ IDs and computed scores. \Cref{fig:5-query-path} distinguishes the
three runtime selection decisions; the ground-truth ordering is an
external evaluation object.

Four quantities must be kept separate:
\begin{align}
 s_i(q)&=\|\mu_i\|_2^2-2\langle q',\mu_i\rangle,
                   &&\text{centroid routing score},\nonumber\\
 \primarydist(q,x)&=\|q'-\widehat p(x)\|_2^2,
                   &&\text{primary score},\nonumber\\
 \refineddist(q,x)&=\|q'-\widehat p(x)-\widehat r(x)\|_2^2
                       +\|q'\|_2^2,
                   &&\text{refined ranking score},\nonumber\\
 \rawdist(q,x)&=\|q-x\|_2^2,
                   &&\text{raw-space distance}.
 \label{eq:5-ranking-objects}
\end{align}
Here $\mu_i$ is an IVF centroid in transformed coordinates,
$\widehat p(x)$ is the primary reconstruction and $\widehat r(x)$ is
the decoded residual. The runtime refined score includes a query-only
constant, as derived from the implemented correction in \Cref{sec:5-6}.
It preserves the ordering of reconstructed distances for a fixed query,
but neither it nor the primary score is the original-vector distance.

\begin{figure}[tbp]
\centering
% Conceptual dataflow; no measured values or GPU timing claims.
\begin{tikzpicture}[font=\small,
 box/.style={draw=ThesisNavy!65,rounded corners=2pt,align=center,
  text width=3.7cm,minimum height=3cm,inner sep=5pt},
 arr/.style={-{Latex[length=2mm]},draw=ThesisNavy!80,thick}]
\node[box,fill=ThesisNavy!5] (route) at (0,0)
 {$q\to q'=\Pi q$\\[3pt]
  Centroid scores $s_i$\\Select $n_{\rm list}\to n_{\rm probe}$ lists\\Visited universe $U$};
\node[box,fill=ThesisNavy!5] (primary) at (4.9,0)
 {Build LUTs\\Read primary codes\\[3pt]
  Score $n_c$ candidates\\Retain primary top-$c$};
\node[box,fill=ThesisNavy!5] (refine) at (9.8,0)
 {Read residual codes\\Compute $\refineddist$\\[3pt]
  Select $c\to k$\\Return IDs and scores};
\draw[arr] (route) -- node[above,font=\footnotesize] {$n_c$} (primary);
\draw[arr] (primary) -- node[above,font=\footnotesize] {$c$} (refine);
\node[draw=ThesisNavy!40,rounded corners=2pt,align=center,
 text width=12.6cm,inner sep=6pt] at (4.9,-2.7)
 {Evaluation only: compare returned IDs with raw-space ground truth under
  $\rawdist$\\Ground truth is outside the compressed search decisions.};
\end{tikzpicture}

\caption[Query stages and selection boundaries]{Query execution and its separate selection boundaries. Routing
selects lists, primary screening selects vector candidates, and final
selection ranks refined survivors. Raw-space ground truth is used for
evaluation; it is not an input to these search decisions.}
\label{fig:5-query-path}
\end{figure}

Routing defines the visited universe $U$ of $n_c$ candidate occurrences;
primary screening retains up to $c=\max\{k,\lceil\alpha k\rceil\}$ for
residual scoring. \Cref{tab:5-query-stages} records each stage's inputs,
reads and costs; \cref{sec:5-2,sec:5-3,sec:5-4,sec:5-5,sec:5-6,sec:5-7}
develop their implementations.

The ordering and exactness columns of \cref{tab:5-query-stages}
attach each claim to its score and tie rule. Chapter 7 tests survivor
sets, ordered outputs and raw-space recall separately.


\begin{table}[tbp]
\centering
\caption[Query-stage selection contracts]{Query-stage contracts. $c_v=\min(c,n_c)$ counts valid primary
survivors. Exactness refers to the stated computed scores; routing and
quantised reconstruction determine raw-space approximation.}
\label{tab:5-selection-contracts}
\label{tab:5-query-stages}
\footnotesize\setlength{\tabcolsep}{3pt}
\begin{tabularx}{\linewidth}{@{}>{\raggedright\arraybackslash}p{1.7cm}>{\raggedright\arraybackslash}p{2.1cm}*{3}{>{\raggedright\arraybackslash}X}@{}}
\toprule
Stage & Input $\to$ output & Data read & Main work & Order / exactness \\
\midrule
Transform & $q\to q'$ & Query; transform & Apply $\Pi$ & Transform (precision in \cref{sec:7-precision}) \\
Probe & $n_{\rm list}\to n_{\rm probe}$ lists & FP32 centroids; list lengths & Score; radix; rank count & $(s_i,i)$; exact probe selection \\
LUT & $q'\to$ tables & Primary codebook & $MD_s2^{8/\glut}$ terms & Real-arithmetic factorisation; FP32 rounding \\
Primary, $\gscan=0$ & $n_c\to c_v$ & Packed primary codes; LUT & $O(n_cM)$ scoring; compaction & Exact distance-order selection; arbitrary cutoff ties \\
Primary, $\gscan\geq1$ & $n_c\to c_v$ & Same primary representation & Score; compact; merge & $(\primarydist,p)$; exact under stated premises \\
Refine & $c_v\to c_v$ scores & Residual codes; correction & $O(c_vd)$ & Computed $\refineddist$; lane-dependent FP32 association \\
Final & $c_v\to k$ & Refined scores; IDs & $O(kc)$ slot scans & $(\refineddist,\mathrm{ID})$; exact with $k$ finite valid entries \\
\bottomrule
\end{tabularx}
\end{table}


\Cref{tab:5-query-stages} gives the stage contracts. The E1/budget
harness explicitly selects 16 refinement lanes and scan policy $-1$
(scheduling among 0/2/4/8), versus library defaults of scalar refinement
and policy 0. Shared bounds are enabled but used only at $\gscan>1$;
the harness calibrates scheduling with 512 database pseudoqueries,
a step the library does not perform automatically. Component
ablations retain their stated overrides; \cref{tab:a-components} maps
mechanisms to recorded implementations.
